\chapter{机器学习究竟在学习什么}

\begin{quote}\small
\textit{本章摘要。}本章定义样本、任务、模型、损失和风险，说明不同损失对应不同预测对象，并讨论正则化、可识别性及预测到决策的边界。
\end{quote}

上一章给出了向量、导数和概率的语言，但能计算一个梯度并不意味着已经定义了值得学习的问题。现在需要把设备观测、未来事件和维护动作分开：观测决定模型能看见什么，标签决定它试图回答什么，损失决定怎样的错误更值得避免。三者不一致时，即使训练完全收敛，也可能得到一个回答错误问题的模型。

\paragraph{本章导语。}
同一段振动记录，可以用来判断轴承是否故障，也可以用来预测下一时刻的振幅。这两个问题需要不同的标签，算错后的后果也不同。本章先把这些差别写清楚，再说明模型中的参数如何通过损失函数学出来。随后讨论正则化、可识别性和部署反馈：为什么训练数据上看起来很好，换一台设备却可能失效？后文介绍各种模型时，都要回到输入、目标和错误代价这几件事。

\section{学习对象：数据、任务与模型}

% auxiliary-resource-guide
本章末的“辅助学习材料”列出与核心方法对应的讲解和实践入口，可在阅读相关推导后，按所列小节对照学习。

\paragraph{从数据到任务。}
数据（data）是对对象、过程或环境的观测记录。一个样本（sample）是一次相对完整的观测，通常记为$\bm{x}\in\R^d$；监督学习中还伴随标签$y$，于是样本写为$(\bm{x},y)$。这里的“样本”不是数据表的一行这么简单：一次设备运行周期、一张图像、一个句子或一段传感器窗口都可以成为样本，关键在于它是否对应一个明确的可学习单位。

“任务”由输入空间$\mathcal{X}$、输出空间$\mathcal{Y}$以及评价准则共同决定。输入空间是模型允许接收的对象集合（例如时序数据、图像、文本等），输出空间是允许回答的结果集合，这两个集合先于可学习参数的确定。分类的输出是离散类别，回归的输出是连续数值，预测性维护可能输出剩余寿命，异常检测则可能没有逐样本标签。相同的数据可以支持不同任务。例如，振动信号既可以用于故障分类，也可以用于生成未来波形或估计设备状态。

\paragraph{把一个任务写成可核查的样本。}
假设每小时读一次温度和振动，把这两项读数排成输入向量：
\[
 \bm x=(\text{温度},\text{振动})^{\mathsf T}\in\R^2.
\]
标签为下一小时温度$y\in\R$时，这就是回归任务；改成“未来一天是否故障”的$y\in\{0,1\}$时，就成为二分类。

监督学习用已知输入和标签学习对应关系，无监督学习则在没有这些目标标签时寻找输入结构，即发现样本之间的相似性、变量之间的关联等规律。例如，不提供故障标签，仅根据温度和振动把相似的设备记录分组，但这些分组不自动代表正常或故障。

同一数据记录能面向不同任务，不同任务同样需要不同的标签或指标。训练是利用已有样本调整参数；推理是在参数训练固定后处理一个新输入，得到对新样本标签或指标的预测。


\paragraph{模型、参数与预测。}
模型根据对数据和标签的观测方式不同，可分为两类：\textbf{确定性模型}和\textbf{概率模型}。确定性模型（deterministic model）是带有可调参数的函数族
\begin{equation}
 f_{\bm{\theta}}:\mathcal{X}\rightarrow\mathcal{Y},
 \qquad \bm{\theta}\in\Theta.
\end{equation}
给定输入$\bm{x}$，确定性模型直接输出$\hat{y}=f_{\bm{\theta}}(\bm{x})$。参数$\bm{\theta}$不是人为凭空指定的数值，而是由训练数据和训练目标共同决定的可学习量。模型结构规定了“允许怎样形式的映射”（线性形式 or 非线性形式），参数则规定了“采用哪一个特定的映射”。通常，这类模型也被称为“\textbf{确定性判别模型}”，因为模型函数拟合的是观测数据到标签的映射。

至于函数族的形式，应是由数据观测$x$与标签$y$（任务）的形式决定的。例如，“图像+分类任务”适用的函数族可以是ResNet、ViT或MLP；“时序+回归任务”适用的函数族可以是RNN、LSTM或TCN。函数族的选择通常由数据类型、任务类型和计算资源共同决定。

对应地，有确定性判别模型，就有“\textbf{概率判别模型}”。概率判别模型输出条件分布$p_{\bm{\theta}}(y\mid\bm{x})$，而不是直接输出$\mathcal Y$中的一个标签。例如，$K$类分类任务的概率输出是$K$个非负且总和为一的数，回归任务的概率输出可以是预测均值和噪声大小（在对标签的噪声建模为高斯的情形下）。确定性输出可以看作概率模型的特殊情形，例如$p(y\mid\bm{x})$集中在单点$f_{\bm{\theta}}(\bm{x})$上。因此，确定性模型与概率模型不是互斥的两个世界，而是表达不确定性程度不同的两种建模方式。

需要注意的是，这里的例子只把随机性放在模型的最终输出端，并不意味着概率模型中的概率性仅存在于输出层。模型内部的隐藏状态可以是随机变量，例如用潜变量$z$表示未直接观测的设备状态，并通过$p_{\bm{\theta}}(y\mid\bm{x},z)$和$p_{\bm{\theta}}(z\mid\bm{x})$共同描述预测过程；模型参数本身也可以被视为随机变量，例如贝叶斯模型用参数分布$p(\bm{\theta}\mid\mathcal D)$表达训练数据给参数带来的不确定性。因而，概率建模可以作用于输出、隐藏变量或参数等不同层次，前面的输出分布只是其中最简单的两种情形。

同时，有一类区别于判别模型的模型，它们叫做“\textbf{生成模型}”。与建模$\bm{x}$到$y$的映射的判别模型不同，生成模型更关注观测$\bm{x}$的生成方式，即建模联合分布$p_{\bm{\theta}}(\bm{x},y)$。生成模型可以用来生成新的样本，也可以用来计算条件分布$p_{\bm{\theta}}(y\mid\bm{x})$，但它们的训练目标和假设通常与判别模型不同。

\section{让模型有点追求}

\paragraph{给定一个模型的训练目标——损失函数。}
训练的最终目的，是在模型允许的参数、隐藏状态或其他可训练变量中，找出使训练目标最优的取值或状态。“最优”并不是脱离任务而自然存在的性质，而是由目标函数规定的；损失函数正是把模型产生的输出与真实任务目标连接起来的媒介。一个好的损失函数应当能够充分反映任务真正关心的目标，并且对任务中的关键差异敏感。例如，分类任务关心的是类别概率及其相对关系，交叉熵通常比直接把“分对或分错”作为训练信号更适合优化。

损失函数的具体形式还取决于标签情况和模型的建模方式。对确定性判别模型，通常先用函数衡量真实标签与模型输出之间的差异，例如回归模型$f_{\bm{\theta}}(\bm{x})$可以采用
\begin{equation}
 \ell_{\mathrm{sq}}\bigl(f_{\bm{\theta}}(\bm{x}),y\bigr)=\frac12\bigl(f_{\bm{\theta}}(\bm{x})-y\bigr)^2,
\end{equation}
再对样本损失求和或取平均。对概率判别模型，若把参数视为随机变量并给定先验$p(\bm{\theta})$，则可以最大化后验
\begin{equation}
 \hat{\bm{\theta}}_{\mathrm{MAP}}
 =\argmax_{\bm{\theta}}p(\bm{\theta}\mid\mathcal D)
 =\argmax_{\bm{\theta}}\bigl[\log p(\mathcal D\mid\bm{\theta})+\log p(\bm{\theta})\bigr].
\end{equation}
等价地，也可以最小化负对数后验；不过在更一般的语境下，这类量通常称为目标函数（objective），不一定称为损失函数。对概率生成模型，则常通过最大化观测数据的对数似然
\begin{equation}
 \hat{\bm{\theta}}_{\mathrm{ML}}
 =\argmax_{\bm{\theta}}\log p_{\bm{\theta}}(\mathcal D)
 =\argmax_{\bm{\theta}}\sum_{i=1}^{n}\log p_{\bm{\theta}}(\bm{x}_i,y_i)
\end{equation}
来确定参数，或等价地最小化负对数似然。因而，损失或目标的构建必须同时考虑任务类型、标签的可观测性以及模型的概率或确定性建模方式。

\paragraph{经验风险与总体风险.}  
有限样本上的拟合与总体上的预测能力需要分开讨论，经验分布到真实分布的一致收敛是统计学习理论的基础问题之一\cite{vapnik1971uniform}。
设$(X,Y)\sim P$描述部署环境中的输入与标签，$X\in\R^d$，$Y$可为实数或类别。总体风险定义为
\begin{equation}
 R(f)=\E_{(X,Y)\sim P}[\ell(f(X),Y)].
\end{equation}
这里的期望包含尚未观察到的未来样本，因此通常不能直接计算。若训练样本独立同分布于$P$（independent and identically distributed，彼此独立且共享同一个分布），对一个事先固定且损失可积的$f$，有$\E[\hat R(f)]=R(f)$。训练前允许选择的模型参数属于参数空间$\Theta$；“事先固定”强调不能先看测试结果再挑选$f$而仍使用这条无偏性解释。若损失方差有限，独立性进一步给出
\begin{equation}
 \Var(\hat R(f))=\frac{1}{n^2}\sum_{i=1}^n\Var(\ell(f(X_i),Y_i))
 =\frac{\Var(\ell(f(X),Y))}{n}.
\end{equation}
这解释了平均损失为什么可用于估计风险，也暴露了限制：相邻设备窗口相关时有额外协方差项；$f$由同一训练集选择后，上述固定模型的无偏论证不能直接用于其训练误差。验证与测试的必要性由此产生，而不是来自惯例。

训练时通常最小化经验风险，真正希望最小化的却是部署分布上的总体风险。“\textbf{过拟合}”的核心表现，是模型在训练集上的经验风险很小，但在未参与训练的新样本上的总体风险仍然较大，因而出现较大的泛化差距。模型容量过大时，它不仅可以学习输入与标签之间可迁移的规律，也可能记住训练样本中的噪声、偶然相关性甚至标签错误；但模型容量过大并不是过拟合的必要条件，训练样本过少、标签噪声较强或反复利用验证集也可能造成过拟合。当训练集有限、包含噪声，或训练集与部署分布存在差异时，这些被模型记住的细节可能无法在新样本中重复出现。典型地，训练早期经验风险和总体风险可能同时下降；训练后期经验风险继续下降，而总体风险可能转而上升，因为模型开始更多地拟合训练数据中的偶然性。


\section{归纳偏置}

前文提到了过拟合，通俗来讲指的是训练误差很低而新样本误差较高。与之对应的一个概念叫做“\textbf{泛化性}”，指模型在“未参与训练和选择的数据”上的有效程度。使用损失函数或目标函数进行训练，直接下降的是经验风险；但我们真正关心的是部署分布上的总体风险。经验风险最小化并不自动意味着总体风险最小化，因此需要在模型构造和训练目标中加入能够约束学习过程的偏好。下面从模型角度和训练目标角度分别说明这种约束如何发挥作用。

正如前文所说，模型本质上是一个函数族；无论这个函数族是确定性的还是含有随机变量，它都不可能平等地偏好所有可能的输入输出关系。模型结构、参数化方式以及训练算法对哪些解更容易被找到、哪些解更容易被保留所作的假设，统称为“\textbf{归纳偏置}”（inductive bias）。归纳偏置的作用，是在训练样本不足以唯一决定规律时，帮助模型从许多能够解释训练数据的候选规律中选择更可能适用于新样本的规律。

例如，对同一图像分类任务，卷积神经网络把局部连接和参数共享写进了模型结构，因此偏好局部模式以及对平移具有一定稳定性的函数；全连接网络则没有这种结构性限制。若图像任务确实具有这些性质，卷积网络就能把有限样本用于学习更有用的变化，而不必分别记住每个像素位置的完全独立关系，因此通常更容易获得较好的泛化性能。这里的“更适合”是统计意义上的，不是保证在任何数据集上都更好。

归纳偏置与可解释性有关，但二者不是同一个概念。可解释性讨论人能否理解模型的预测依据；归纳偏置讨论模型在学习时预先接受了哪些结构假设。一个模型可以具有明确的归纳偏置，却仍然难以解释；也可以相对容易解释，但其归纳偏置并不适合当前任务。判断归纳偏置是否良好，关键在于它是否与数据生成过程和任务目标相匹配，并能否在有限样本下缩小训练风险与总体风险之间的差距。

正则化（regularization）提供了加入这种偏好的另一条途径：不一定更换模型结构，而是在训练目标中明确惩罚某些不希望出现的解。其基本形式是
\begin{equation}
 \min_{\bm{\theta}}\;\hat{R}(\bm{\theta})+\lambda\Omega(\bm{\theta}),
 \qquad \Omega(\bm{\theta})=\frac12\|\bm{\theta}\|_2^2.\label{eq:foundations-regularized-risk}
\end{equation}
这里第一项$\hat R(\bm{\theta})$要求模型解释训练数据，第二项$\Omega(\bm{\theta})$表达我们对模型复杂度的偏好。选择$\Omega(\bm{\theta})=\frac12\|\bm{\theta}\|_2^2$，表示在其他条件相近时偏好范数较小的参数；它并不声称真实参数一定接近零，而是把“大权重可能对应不稳定拟合”作为一种可调的建模偏好。系数$\lambda\geq0$决定这种偏好有多强：$\lambda=0$时只拟合训练数据；$\lambda$增大时，模型愿意牺牲一部分训练集拟合精度，以换取更小的参数和可能更好的泛化。正则化的目的不是让训练误差尽可能低，而是在经验风险和对复杂度的约束之间寻找适合部署的平衡。

\begin{example}
 对于温度传感器预测下一时刻温度，模型并不需要学习所有历史序列到未来数值的任意映射。若已知相邻时刻通常连续变化，可以采用平滑性偏置；若设备存在周期，则可以把周期特征加入输入。模型的“聪明”来自数据，也来自这些合理的限制。
\end{example}

\paragraph{惩罚形式与约束形式如何联系。}
为什么可以把正则化理解成“控制复杂度”？一种直接的写法是先规定允许的复杂度上限：
\begin{equation}
 \min_{\theta}\hat R(\theta)\quad\text{subject to}\quad \Omega(\theta)\leq c.
\end{equation}
这句话的意思是：在所有不超过复杂度预算$c$的模型中，选择训练风险最小的模型。它明确表达了两件事：训练数据需要拟合，但模型不能为了拟合少量样本而无限增加复杂度。

约束问题也可以通过拉格朗日乘子写成
\begin{equation}
 L(\theta,\lambda)=\hat R(\theta)+\lambda\bigl(\Omega(\theta)-c\bigr),\qquad \lambda\geq0.
\end{equation}
固定$\lambda$后，$-\lambda c$与$\theta$无关，不会改变参数最优解，于是得到惩罚形式$\hat R(\theta)+\lambda\Omega(\theta)$。因此，$c$和$\lambda$分别是两种不同的控制方式：前者规定复杂度上限，后者规定超出较简单解的代价。在线性凸问题并满足适当严格可行性条件时，二者可以通过拉格朗日最优性条件联系起来；对非凸网络，不能据此声称任意约束问题都与某个惩罚问题全局等价，也不能认为$\lambda$与$c$必然一一对应。

在约束问题的最优点$\hat\theta$处，若目标与约束满足相应的凸性条件，可以写出
\begin{equation}
 0\in\partial\hat R(\hat\theta)+\lambda\partial\Omega(\hat\theta),\qquad
 \Omega(\hat\theta)\leq c,\qquad
 \lambda(\Omega(\hat\theta)-c)=0.
\end{equation}
$\partial$是允许不可微凸函数的次梯度集合。最后一个等式是互补条件。若$\Omega(\hat\theta)<c$，说明复杂度预算没有用完，限制没有起作用，此时可以有$\lambda=0$；若$\Omega(\hat\theta)=c$，说明模型已经碰到复杂度边界，进一步降低训练风险通常需要增加复杂度，约束因此会通过正的$\lambda$影响最优解。这里的“补偿”不是说两个数必须相等，而是说训练风险的改善必须值得为增加复杂度所付出的代价。

最后用一维例子看清“收缩”究竟发生了什么。设无正则化时的训练风险为
\begin{equation}
 \hat R(w)=\frac12(w-a)^2,
\end{equation}
其中$a$是只根据训练数据得到的最优估计。它的最小值在$w=a$处，因为此时$w$与$a$完全一致。现在加入平方惩罚，训练目标变为
\begin{equation}
 J(w)=\frac12(w-a)^2+\frac{\lambda}{2}w^2.
\end{equation}
第一项希望$w$接近数据给出的估计$a$，第二项希望$w$接近零。对$w$求导得到
\begin{equation}
 J'(w)=(w-a)+\lambda w.
\end{equation}
令导数为零，得到$(1+\lambda)w=a$，因此
\begin{equation}
 \hat w=\frac{a}{1+\lambda}.
\end{equation}
由于$\lambda\geq0$，有$|\hat w|\leq|a|$；当$\lambda=0$时不收缩，随着$\lambda$增大，$\hat w$逐渐靠近零。这就是“收缩”的含义：正则化不会自动判断某个变量是噪声，也不会直接把所有不重要的系数删掉，而是把参数估计向预先选定的中心（这里是零）拉近。

收缩为什么可能改善总体风险？训练数据中的偶然波动会使无正则化估计$a$在不同训练集之间变化很大，这种变化称为估计的不稳定性或方差。平方惩罚牺牲了一部分对当前训练集的贴合，使估计更稳定；但如果真实系数确实远离零，收缩会带来偏差。因此，正则化是在偏差与方差之间作权衡：只有方差的降低足以抵消偏差增加，使用验证数据选择的$\lambda$才可能降低总体预测风险。它不是“惩罚越强越好”，也不是“参数越小越接近真实机制”。

模型选择至少同时受四类约束：数据规模与标注质量、任务中的依赖结构、训练与推理资源、错误的业务代价。复杂模型并不自动优于简单模型；它只有在表示了简单模型无法表示的结构时才有必要。后文讨论的序列、视觉和语言模型，实际上是在不同结构约束下对这四个问题的回答。


\section{任务的可学习性}
在选择模型之前，还应先问一个更根本的问题：给定的输入中，是否含有足以预测目标的稳定信息？例如，设输入是风电机组叶片的图像，标签是同一时刻的发电功率。凭借领域知识可以预期，功率还强烈依赖风速、风向、空气密度、机组运行状态和控制策略；一张只反映叶片外观的图像，通常不足以充分决定功率。图像中也许包含结冰、污损或损伤等与功率有关的间接线索，但这种联系必须由数据和机理共同证实，不能因为两个变量出现在同一张数据表中，就假定二者必然存在可用于预测的关系。

训练集上的拟合能力不能单独回答“任务是否可以被学习”这个问题。只要模型足够灵活，它就可以把输入空间划分为许多很小的区域，并为每个训练样本建立对应映射。以决策树回归为例，当训练输入彼此可区分，且树的深度、叶节点样本数等不受限制时，树可以把训练样本逐个分到不同叶节点，从而达到零训练误差；分类树也可以达到百分之百训练准确率。这说明模型有能力记住训练样本，却不说明图像中存在对未来功率稳定有效的预测线索。若新图像落入训练时未被充分覆盖的区域，模型只能依赖未经验证的划分和映射，训练集上的完美表现便可能无法迁移到新样本。这正是过拟合的一种典型表现。

其实，这个问题很简单，只需要判断一件事：“相似的”输入是否具有“相似的”标签。从统计角度看，一个监督任务在给定输入$X$时是否可学习，取决于真实条件分布$P(Y\mid X)$是否具有可利用的结构，以及这一结构能否在训练环境和部署环境之间保持足够稳定。这里的$P(Y\mid X)$描述“在给定输入时，标签可能如何变化”。若能够在输入空间的各个相关区域持续获得带标签样本，并且条件分布在邻近输入之间不会发生任意剧烈的变化，那么观察到的局部模式才有可能外推到附近的未见样本。连续性、平滑性或低维结构都可以成为这种外推的依据。那么，在足够的样本覆盖、合适的模型族的条件下，任务的训练才具备泛化性。

现实中的困难在于，输入空间往往高维，且样本真实分布不均匀；有限样本更不可能把整个空间都覆盖。因此，数据中是否存在规律通常不能仅靠观察几个样本来断言。一个有用的直觉是：若相同或非常相近的输入在输入空间的一个小邻域内反复对应差异很大的标签，那么只使用这些输入时就存在不可消除的不确定性；即使知道真实条件分布，也不可能准确预测每一个样本。相反，若$X$能够稳定地缩小$Y$的不确定性，使合适的预测函数在独立新样本上显著优于“不看输入、只预测平均值或多数类”的基线，则该任务具有可学习的信号。这里的“规律”不要求是线性的、显式的，或容易写成公式；它可以是高度非线性的统计关系，但必须能够在新的数据中重复出现。

有限数据无法从逻辑上证明这种规律一定存在，只能不断积累证据。实践中应结合四类检查：首先用领域机理审查输入到标签的路径，确认不是把结果发生之后的信息误当作输入；其次检查标签的测量质量、时间对齐和样本覆盖，避免真实信号被噪声或错误配对淹没；再次在严格隔离的验证集或测试集上，与不使用输入的简单基线比较；最后检验结果是否能跨时间、机组或运行条件复现。若模型只在训练集上优于基线，或一旦按时间、设备划分测试集就失效，更合理的结论通常是现有输入不足、分布发生了变化，或模型只记住了偶然性，而不是急于增加模型复杂度。

\paragraph{从训练到部署及反馈。}
沿图\ref{fig:foundations-learning-loop}从采集到上线的路径，可以依次看四个阶段。先把设备记录整理成输入和标签，再选择能处理这些输入的模型，通过训练求出参数，最后用这些参数预测新记录。上线后的错误又会提示我们补数据、改标签或重新训练。评价时除了训练损失，还要检查新工况下的表现、概率是否校准、计算是否及时，以及误报和漏报的后果。

\begin{figure}[htbp]
\centering
\begin{tikzpicture}[>=Latex, node distance=0.85cm, every node/.style={font=\small}]
 \node[draw, rounded corners, fill=blue!10, minimum width=2.15cm, minimum height=0.9cm] (data) {\shortstack{数据\\$\mathcal D$}};
 \node[draw, rounded corners, fill=orange!12, right=of data, minimum width=2.35cm, minimum height=0.9cm] (task) {\shortstack{任务定义\\$X,Y,\ell$}};
 \node[draw, rounded corners, fill=green!12, right=of task, minimum width=2.35cm, minimum height=0.9cm] (model) {\shortstack{模型\\$f_\theta$}};
 \node[draw, rounded corners, fill=yellow!20, right=of model, minimum width=2.35cm, minimum height=0.9cm] (train) {\shortstack{训练\\$\argmin_\theta$}};
 \node[draw, rounded corners, fill=red!10, right=of train, minimum width=2.15cm, minimum height=0.9cm] (deploy) {\shortstack{推理\\与部署}};
 \draw[->, thick] (data)--(task); \draw[->, thick] (task)--(model); \draw[->, thick] (model)--(train); \draw[->, thick] (train)--(deploy);
 \draw[->, thick, dashed] (deploy.north) .. controls +(0,0.72) and +(0,0.72) .. node[above,font=\scriptsize]{反馈、漂移、错误分析} (data.north);
\end{tikzpicture}
\caption{工业数据智能的训练、部署与反馈：部署反馈会重新影响数据和任务定义。}
\label{fig:foundations-learning-loop}
\end{figure}

\section{尚未解决的问题——优化}

本章讲解了数据、任务、目标函数、泛化性等概念，告诉了大家什么是一个好的模型，应该用怎样的目标函数训练一个模型，应该怎样评价一个模型。但是还缺一点，有了目标函数，模型中的可学习参数（或者是先验分布），如何根据目标函数一步一步被训练到“最优”状态。关于这个问题，我们后面会用一个单独的章节去介绍。
\section*{辅助学习材料}
\begingroup
\emergencystretch=3em
\begin{itemize}
\item \textbf{Linear Models}；作者/平台：scikit-learn Developers；英文；官方文档。重点阅读 Ordinary Least Squares、Ridge、Lasso 与 Logistic Regression 小节，配合本章的线性预测、正则化和分类损失理解参数约束与风险最小化。\href{https://scikit-learn.org/stable/modules/linear\_model.html}{在线阅读}。
\item \textbf{Cross-validation: evaluating estimator performance}；作者/平台：scikit-learn Developers；英文；官方文档。重点阅读训练/验证/测试划分、交叉验证和数据泄漏说明，帮助把本章的经验风险、部署风险与可复现实验边界落实为流程。\href{https://scikit-learn.org/stable/modules/cross\_validation.html}{在线阅读}。
\item \textbf{损失函数（Loss Function）}；知乎；中文解说。比较平方损失与交叉熵，核对本章预测对象；标题与主题据必应索引，原页受限。\href{https://zhuanlan.zhihu.com/p/261059231}{在线阅读}。
\item \textbf{损失函数（lossfunction）的全面介绍（简单易懂版）}；CSDN；中文博客。阅读 L1、MSE 与交叉熵说明，比较不同误差的代价；已核对公开页面标题与索引摘要。\href{https://blog.csdn.net/weixin\_57643648/article/details/122704657}{在线阅读}。
\item \textbf{2021 - 类神经网络训练不起来怎么办(四) 损失函数 (Loss) 也可能有影响}；李宏毅；bilibili；中文课程。选看第22集，对照本章损失与学习目标。2026年10月4日官方公开元数据显示整套课程累计播放约299.5万，并非本分集播放量；已核对目录与计数，未逐帧观看。\href{https://www.bilibili.com/video/BV1Wv411h7kN/?p=22}{观看分集}。
\end{itemize}
\endgroup

